% Generated from locale/tr/content/set-theory/z/arbintersections.tex; do not hand-edit.


\olfileid[tr]{sth}{z}{arbintersections}
\olsection{Ek: Kapanış, Küme Oluşturma ve Kesişim}

\olref[sth][z][milestone]{sec} içinde, \olref[sfr][][]{part} kısmındaki
naif çalışmaya geri dönüp bunun $\Zminus$ içinde yürütülebileceğini
denetlemenizi önermiştik. Bu öğüdü izlediyseniz \emph{bir} nokta sizi
zorlamış olabilir: Dedekind'in \emph{kapanışları} ele alışında
\emph{kesişim} kullanması.

\olref[sfr][infinite][dedekind]{Closure} içinden şunu hatırlayın:
\begin{align*}
  \closureofunder{f}{o} & = \bigcap\Setabs{X}{o \in X \text{ ve $X$,
$f$ altında kapalıdır}}.
\intertext{Bunun genel biçimi şu türden bir tanımdır:}
  C & = \bigcap\Setabs{X}{\phi(X)}.
\end{align*}
Ancak bu tehlike çanlarını çaldırmalıdır: Naif Küme Oluşturma başarısız
olduğundan $\Setabs{X}{\phi(X)}$ kümesinin var olduğunun güvencesi yoktur.
Öyleyse bu tür tanımların \emph{hile yaptığı} kuşkusunu doğurması doğaldır.

Neyse ki hile yapmıyorlar; daha doğrusu, oldukları biçimiyle hile
yapıyorlarsa onları meşru kılmak için dürüstçe çalışabiliriz. Bu dürüst
çalışma, her $A\neq\emptyset$ için $\bigcap A$'nın neden var olduğunu
açıkladığımız \olref[sth][z][sep]{prop:intersectionsexist} içinde önceden
haber verilmişti. Şimdi bunu açıkça ortaya koyacağız.

Kaplamsallık verildiğinde, $C$'yi $\bigcap\Setabs{X}{\phi(X)}$ olarak
tanımlamaya çalışırken aslında yalnızca şu koşula uyan bir $C$ nesnesi
istiyoruz:
\begin{equation}\ollabel{bicondelimarbintersection}
	\forall x(x \in C \liff \forall X(\phi(X) \lif x \in X))\tag{*}
\end{equation}
Şimdi $\phi(S)$ koşulunu sağlayan \emph{bir} $S$ kümesinin bulunduğunu
varsayın. \olref{bicondelimarbintersection} koşulunu sağlamak için $C$'yi
Ayırma'yla şöyle tanımlayabiliriz:
\[
	C = \Setabs{x \in S}{\forall X(\phi(X) \lif x \in X)}.
\]
Bu tanımın istendiği gibi \olref{bicondelimarbintersection} koşulunu
verdiğini denetlemeyi alıştırma olarak bırakıyoruz.
%  fix $x$. First suppose that $x \in C$, as (re)defined; then both $x
%  \in S$ and $\forall X(\phi(X) \lif x \in X)$; but since
%  $\phi(S)$, this reduces to the condition that $\forall X(\phi(X)
%  \lif x \in X)$. Conversely, suppose $\forall X(\phi(X) \lif
%  x \in X)$; then, since $\phi(S)$, we also have that $x \in S$; so
%  $x \in c$, as (re)defined. 
Bu genel strateji, kesişimleri tanımlarken Naif Küme Oluşturma'nın görünürdeki
her kullanımından kaçınmamızı sağlar. Bizi bu düşünce çizgisine getiren özel
durumda, yani $\closureofunder{f}{o}$ için bunun nasıl işleyeceğini görelim.
\olref[sfr][infinite][dedekind]{closureproperties} ispatına,
$o\in\ran{f}\cup\{o\}$ ve $\ran{f}\cup\{o\}$ kümesinin $f$ altında kapalı
olduğunu belirterek başlamıştık. Dolayısıyla istediğimiz şeyi şöyle
tanımlayabiliriz:
\[
	\closureofunder{f}{o} = \Setabs{x \in \ran{f} \cup \{o\}}{(\forall X \ni o)(X \text{, $f$ altında kapalıdır} \lif x \in X)}.
\]

